\subsection{Circularity Checks}
\label{app:circularity}

\noindent COCO \texttt{train2017} NPMI defines the paired splits, the STRANGE-Bench tiers and the AAS tilt, and the AAS late stage trains on scenes from the generator of the benchmark. We test each of these links against a source the benchmark was not built from.

\noindent\textbf{Tiers under a second co-occurrence source.} Tab.~\ref{tab:app_circ_tiers} scores every anchor-partner pair of STRANGE-Bench with the NPMI of Objects365 \texttt{train} ($1.70$M images mapped to the COCO classes), which played no part in building the benchmark. Objects365 ranks the typical above the atypical pairs at every level (AUC $0.97$ to $0.99$, against $0.99$ to $1.00$ for COCO). If classes co-occurred independently, $85\%$ of the pairs that never co-occur in COCO \texttt{train2017} would be expected there at least $10$ times, where observing none has probability below $5{\times}10^{-5}$, so their absence is not a sampling accident. The pairs of $T_{\text{never}}$ are rare in Objects365 as well: $36.4\%$ of them never co-occur there, and their median Objects365 NPMI is $-0.22$.

\begin{table*}[!htp]%
\centering
\footnotesize
\setlength{\tabcolsep}{4pt}
\renewcommand{\arraystretch}{0.9}
\caption{\textbf{STRANGE-Bench tiers under Objects365 NPMI.} Objects365, which played no part in building the benchmark, ranks the typical above the atypical pairs at every level (AUC $0.97$ to $0.99$), and $36.4\%$ of the $T_{\text{never}}$ pairs never co-occur there either. Anchor-partner pairs of every scene scored with COCO \texttt{train2017} NPMI, which defines the tiers, and with Objects365 \texttt{train} NPMI; a pair that never co-occurs in a source is scored $-1$. NPMI: median over the typical / atypical arm of a level; $N_{ij}{=}0$: share of atypical pairs that never co-occur in the source; AUC: probability that a typical pair outranks an atypical pair of the same level.}
\label{tab:app_circ_tiers}
\begin{tabular}{l ccc ccc}
\toprule
& \multicolumn{3}{c}{COCO \texttt{train2017} (defines the tiers)} & \multicolumn{3}{c}{Objects365 \texttt{train}} \\
\cmidrule(lr){2-4} \cmidrule(lr){5-7}
Level & NPMI typ.\ / atyp. & $N_{ij}{=}0$ atyp.\ (\%) & AUC & NPMI typ.\ / atyp. & $N_{ij}{=}0$ atyp.\ (\%) & AUC \\
\midrule
$k=1$ & $0.20$ / $-1.00$ & $57.5$ & $1.00$ & $0.18$ / $-0.24$ & $32.2$ & $0.99$ \\
$k=2$ & $0.18$ / $-1.00$ & $53.1$ & $1.00$ & $0.15$ / $-0.21$ & $29.6$ & $0.99$ \\
$k=3$ & $0.18$ / $-0.38$ & $45.9$ & $1.00$ & $0.14$ / $-0.21$ & $28.2$ & $0.98$ \\
$k=4$ & $0.17$ / $-0.36$ & $43.1$ & $0.99$ & $0.14$ / $-0.20$ & $27.2$ & $0.97$ \\
\midrule
$T_{\text{never}}$ & -- / $-1.00$ & $82.4$ & -- & -- / $-0.22$ & $36.4$ & -- \\
\bottomrule
\end{tabular}

\end{table*}

\noindent\textbf{Paired splits defined by Objects365 NPMI.} Re-slicing COCO \texttt{val} with Objects365 NPMI (same construction and image pool, different images) leaves a typicality gap in every row we re-sliced, the baseline, paste and AAS rows of Tabs.~\ref{tab:codino_main} to~\ref{tab:dinov3_main} (Tab.~\ref{tab:app_circ_reslice}). The gap is smaller than under COCO NPMI, by $0.5$ to $3.9$~p.p.\@, so part of the gap of the COCO splits reflects the signal that defines them, but in no row does it disappear.

\begin{table*}[!htp]%
\centering
\footnotesize
\setlength{\tabcolsep}{4pt}
\renewcommand{\arraystretch}{0.9}
\caption{\textbf{COCO \texttt{val} paired splits defined by COCO or by Objects365 NPMI.} Every row keeps a gap on the Objects365-NPMI split, $0.5$ to $3.9$~p.p.\@ smaller than on the COCO-NPMI split. Baseline, paste and AAS rows of Tabs.~\ref{tab:codino_main} to~\ref{tab:dinov3_main} (same checkpoints). Both splits use the same construction and image pool (App.~\ref{app:splits-impl}) but select different images. Detectors: full-annotation AP; DINOv3: F1$_{\text{micro}}$.}
\label{tab:app_circ_reslice}
\begin{tabular}{l ccc ccc}
\toprule
& \multicolumn{3}{c}{Split by COCO NPMI} & \multicolumn{3}{c}{Split by Objects365 NPMI} \\
\cmidrule(lr){2-4} \cmidrule(lr){5-7}
Augmentation & Typ.\,$\uparrow$ & Atyp.\,$\uparrow$ & $\Delta^{\text{T-A}}\downarrow$ & Typ.\,$\uparrow$ & Atyp.\,$\uparrow$ & $\Delta^{\text{T-A}}\downarrow$ \\
\midrule
\multicolumn{7}{l}{\emph{Co-DINO R-50 (AP)}} \\
No aug & 44.2 & 35.5 & \underline{8.7} & 44.3 & 36.2 & 8.2 \\
Std. & \underline{53.2} & 44.0 & 9.2 & \textbf{52.3} & 44.2 & 8.1 \\
SCP & \underline{53.2} & \underline{44.2} & 9.0 & 51.5 & \underline{44.4} & \underline{7.1} \\
XPaste & 52.0 & 43.0 & 9.0 & 50.7 & 43.1 & 7.6 \\
\rowcolor{oursrow} SCP + AAS \emph{(ours)} & \textbf{54.0} & \textbf{45.5} & \textbf{8.5} & \underline{52.2} & \textbf{45.9} & \textbf{6.3} \\
\rowcolor{oursrow} \textbf{XPaste + AAS \emph{(ours)}} & 52.1 & 43.3 & 8.8 & 50.9 & 43.8 & \underline{7.1} \\
\midrule
\multicolumn{7}{l}{\emph{Grounding DINO Swin-B (AP)}} \\
Zero Shot & 57.7 & 51.4 & 6.3 & 56.3 & 50.9 & 5.4 \\
SCP & 59.6 & 53.4 & 6.1 & 58.3 & 53.2 & 5.1 \\
XPaste & \underline{59.7} & \textbf{54.1} & \textbf{5.6} & \textbf{58.6} & \textbf{53.6} & \underline{4.9} \\
\rowcolor{oursrow} SCP + AAS \emph{(ours)} & \textbf{59.8} & 53.6 & 6.2 & \underline{58.5} & \underline{53.3} & 5.1 \\
\rowcolor{oursrow} \textbf{XPaste + AAS \emph{(ours)}} & \underline{59.7} & \underline{53.7} & \underline{6.0} & 58.4 & \textbf{53.6} & \textbf{4.8} \\
\midrule
\multicolumn{7}{l}{\emph{DINOv3 ViT-B/16 (F1$_{\text{micro}}$)}} \\
No aug & 77.8 & 70.1 & 7.7 & 76.7 & 72.9 & 3.8 \\
Std. & 79.8 & 74.3 & 5.5 & 79.8 & 76.1 & 3.7 \\
SCP & 78.8 & 75.1 & \textbf{3.7} & 79.3 & 76.5 & 2.8 \\
XPaste & \underline{80.2} & \underline{76.1} & 4.1 & \underline{80.0} & \underline{77.8} & \textbf{2.2} \\
\rowcolor{oursrow} SCP + AAS \emph{(ours)} & 79.0 & 75.0 & 4.0 & 79.1 & 76.8 & \underline{2.3} \\
\rowcolor{oursrow} \textbf{XPaste + AAS \emph{(ours)}} & \textbf{80.6} & \textbf{76.8} & \underline{3.9} & \textbf{80.6} & \textbf{78.0} & 2.5 \\
\bottomrule
\end{tabular}

\end{table*}

\noindent\textbf{AAS driven by Objects365 NPMI.} We re-ran the XPaste + AAS recipe with Objects365 NPMI wherever NPMI enters training: the generated training scenes, their typical-partner filter and the AAS tilt; the benchmark stays on COCO NPMI. The two sources admit largely the same training pairs ($542$ shared of $603$ COCO and $562$ Objects365 pairs), so this check asks whether the result survives the change of source, not whether it transfers to a disjoint set of pairs. In both families the two AAS rows differ by at most $0.3$~p.p.\@ on any column, about the spread over training draws, and both stay ahead of XPaste on STRANGE-Bench and on its held-out classes (Tab.~\ref{tab:app_circ_o365_aas}).

\begin{table*}[!htp]%
\centering
\footnotesize
\setlength{\tabcolsep}{4pt}
\renewcommand{\arraystretch}{0.9}
\caption{\textbf{XPaste + AAS driven by COCO or by Objects365 NPMI.} The two AAS rows differ by at most $0.3$~p.p.\@ on any column, and both stay ahead of XPaste on STRANGE-Bench and on its held-out classes. Objects365 NPMI replaces COCO NPMI in the generated training scenes, their filter and the AAS tilt; the benchmark stays on COCO NPMI. Means over $n$ training draws; Co-DINO rows are late-stage models trained on mini-COCO (App.~\ref{app:beta_sweep}). Held-out classes: the subset of STRANGE-Bench whose classes never appear in the generated training scenes.}
\label{tab:app_circ_o365_aas}
\begin{tabular}{l c cc cc cc}
\toprule
& & \multicolumn{2}{c}{COCO} & \multicolumn{2}{c}{STRANGE-Bench} & \multicolumn{2}{c}{Held-out classes} \\
\cmidrule(lr){3-4} \cmidrule(lr){5-6} \cmidrule(lr){7-8}
Augmentation & $n$ & val\,$\uparrow$ & $\Delta^{\text{T-A}}\downarrow$ & $\Delta_{1\text{-}5}\downarrow$ & $T_{\text{never}}\uparrow$ & $\Delta_{1\text{-}5}\downarrow$ & $T_{\text{never}}\uparrow$ \\
\midrule
\multicolumn{8}{l}{\emph{Co-DINO R-50, mini-COCO late-stage protocol (AP)}} \\
XPaste & 6 & \textbf{43.3} & \underline{8.0} & 1.0 & \underline{87.3} & \underline{1.4} & 87.7 \\
\rowcolor{oursrow} \textbf{XPaste + AAS, COCO NPMI \emph{(ours)}} & 4 & \underline{43.2} & \textbf{7.8} & \underline{0.2} & \textbf{89.3} & \textbf{0.5} & \underline{89.6} \\
\rowcolor{oursrow} \textbf{XPaste + AAS, Objects365 NPMI \emph{(ours)}} & 2 & \underline{43.2} & 8.1 & \textbf{0.1} & \textbf{89.3} & \textbf{0.5} & \textbf{89.9} \\
\midrule
\multicolumn{8}{l}{\emph{DINOv3 ViT-B/16 (F1$_{\text{micro}}$)}} \\
XPaste & 4 & 82.3 & \underline{4.2} & \underline{$-0.6$} & 95.1 & $-0.9$ & 95.2 \\
\rowcolor{oursrow} \textbf{XPaste + AAS, COCO NPMI \emph{(ours)}} & 4 & \textbf{82.5} & \textbf{4.1} & \textbf{\boldmath $-0.8$} & \textbf{96.1} & \textbf{\boldmath $-1.2$} & \textbf{96.0} \\
\rowcolor{oursrow} \textbf{XPaste + AAS, Objects365 NPMI \emph{(ours)}} & 4 & \underline{82.4} & \underline{4.2} & \textbf{\boldmath $-0.8$} & \underline{95.9} & \underline{$-1.1$} & \underline{95.7} \\
\bottomrule
\end{tabular}

\end{table*}

\noindent\textbf{Held-out classes.} The generated training scenes of the AAS late stage use only part of the classes; the scenes of STRANGE-Bench whose classes all come from the other part form a held-out subset ($3{,}234$ images). On it (Tab.~\ref{tab:app_circ_heldout}), XPaste\,+\,AAS keeps the highest atypical score and $T_{\text{never}}$ in every family, and each AAS row scores higher on both than the row it starts from, with a gap no larger (DINOv3 XPaste and XPaste\,+\,AAS tie at $-1.1$). The gains of the AAS rows therefore do not come from classes seen in generated training scenes.

\begin{table*}[!htp]%
\centering
\footnotesize
\setlength{\tabcolsep}{4pt}
\renewcommand{\arraystretch}{0.9}
\caption{\textbf{STRANGE-Bench and its held-out-class subset.} On the held-out classes, XPaste\,+\,AAS keeps the highest atypical score and $T_{\text{never}}$ in every family, and each AAS row improves both over the row it starts from. Baseline, paste and AAS rows of Tabs.~\ref{tab:codino_main} to~\ref{tab:dinov3_main}; the held-out subset contains only classes that never appear in the generated training scenes of the AAS rows. Means over $T_1$ to $T_5$; detectors: AP; DINOv3: F1$_{\text{micro}}$.}
\label{tab:app_circ_heldout}
\begin{tabular}{l ccc ccc}
\toprule
& \multicolumn{3}{c}{STRANGE-Bench} & \multicolumn{3}{c}{Held-out classes} \\
\cmidrule(lr){2-4} \cmidrule(lr){5-7}
Augmentation & $T_{1\text{-}5}^{\text{A}}\uparrow$ & $\Delta_{1\text{-}5}\downarrow$ & $T_{\text{never}}\uparrow$ & $T_{1\text{-}5}^{\text{A}}\uparrow$ & $\Delta_{1\text{-}5}\downarrow$ & $T_{\text{never}}\uparrow$ \\
\midrule
\multicolumn{7}{l}{\emph{Co-DINO R-50 (AP)}} \\
No aug & 76.4 & 7.0 & 76.3 & 76.6 & 7.3 & 76.7 \\
Std. & 81.9 & 6.3 & 83.3 & 81.6 & 7.3 & 84.1 \\
SCP & 86.3 & 3.9 & 87.5 & 86.1 & 4.9 & 88.1 \\
XPaste & 89.0 & \underline{1.7} & \underline{90.4} & \underline{89.1} & \underline{2.2} & 89.7 \\
\rowcolor{oursrow} SCP + AAS \emph{(ours)} & \underline{89.2} & 2.1 & 90.1 & \underline{89.1} & 2.8 & \underline{90.5} \\
\rowcolor{oursrow} \textbf{XPaste + AAS \emph{(ours)}} & \textbf{90.7} & \textbf{1.0} & \textbf{92.1} & \textbf{90.8} & \textbf{0.9} & \textbf{92.7} \\
\midrule
\multicolumn{7}{l}{\emph{Grounding DINO Swin-B (AP)}} \\
Zero Shot & 90.9 & 2.4 & 92.4 & 91.2 & 2.8 & 93.0 \\
SCP & 92.4 & 0.9 & 93.3 & 92.7 & 1.4 & 93.6 \\
XPaste & 93.6 & \textbf{0.3} & \underline{94.6} & 93.3 & 1.5 & 95.0 \\
\rowcolor{oursrow} SCP + AAS \emph{(ours)} & \underline{93.7} & \underline{0.6} & 94.5 & \underline{93.6} & \underline{1.3} & \underline{95.1} \\
\rowcolor{oursrow} \textbf{XPaste + AAS \emph{(ours)}} & \textbf{94.7} & \textbf{0.3} & \textbf{95.2} & \textbf{94.4} & \textbf{1.1} & \textbf{95.7} \\
\midrule
\multicolumn{7}{l}{\emph{DINOv3 ViT-B/16 (F1$_{\text{micro}}$)}} \\
No aug & 66.3 & 19.2 & 61.5 & 67.4 & 18.9 & 62.7 \\
Std. & 67.5 & 22.2 & 62.6 & 69.0 & 20.9 & 62.8 \\
SCP & 87.1 & 4.3 & 86.5 & 86.8 & 4.5 & 86.3 \\
XPaste & \underline{94.4} & \underline{$-0.7$} & \underline{95.0} & \underline{94.6} & \textbf{\boldmath $-1.1$} & \underline{95.3} \\
\rowcolor{oursrow} SCP + AAS \emph{(ours)} & 91.2 & 1.2 & 89.5 & 90.2 & \underline{2.1} & 88.3 \\
\rowcolor{oursrow} \textbf{XPaste + AAS \emph{(ours)}} & \textbf{95.5} & \textbf{\boldmath $-0.8$} & \textbf{96.1} & \textbf{95.5} & \textbf{\boldmath $-1.1$} & \textbf{95.9} \\
\bottomrule
\end{tabular}

\end{table*}

\noindent\textbf{A second generator.} We re-rendered $999$ held-out-class scenes with SD 3.5 Medium from the captions of the benchmark and passed them through the same audit and filter; enough complete scenes survive for $T_4$ ($67$ per arm) and $T_{\text{never}}$ ($136$). The atypical arm of $T_4$ scores lower on these renders for every model, but every AAS row keeps a smaller $T_4$ gap and a higher $T_{\text{never}}$ than its control (Tab.~\ref{tab:app_circ_xgen}), so the gain of AAS is not tied to the generator of its training scenes. These paste rows are late-stage models of our development runs, trained on mini-COCO (App.~\ref{app:beta_sweep}; recipes in App.~\ref{app:exp-recipes}).

\begin{table*}[!htp]%
\centering
\footnotesize
\setlength{\tabcolsep}{4pt}
\renewcommand{\arraystretch}{0.9}
\caption{\textbf{Held-out-class scenes rendered by the generator of the benchmark (SD 3.5 Large) and by SD 3.5 Medium.} On the SD 3.5 Medium renders the atypical arm of $T_4$ scores lower for every model, yet every AAS row keeps a smaller $T_4$ gap and a higher $T_{\text{never}}$ than its control. Both generators render from the same captions, at $k = 3$ ($T_4$) and for $T_{\text{never}}$. Paste rows: mini-COCO late-stage models, means over training draws.}
\label{tab:app_circ_xgen}
\begin{tabular}{l cccc cccc}
\toprule
& \multicolumn{4}{c}{SD 3.5 Large (benchmark generator)} & \multicolumn{4}{c}{SD 3.5 Medium} \\
\cmidrule(lr){2-5} \cmidrule(lr){6-9}
Augmentation & $T_4^{\text{T}}\uparrow$ & $T_4^{\text{A}}\uparrow$ & $\Delta_4\downarrow$ & $T_{\text{never}}\uparrow$ & $T_4^{\text{T}}\uparrow$ & $T_4^{\text{A}}\uparrow$ & $\Delta_4\downarrow$ & $T_{\text{never}}\uparrow$ \\
\midrule
\multicolumn{9}{l}{\emph{Co-DINO R-50 (AP; paste rows: mini-COCO late-stage protocol)}} \\
Std. & \underline{87.3} & 79.7 & 7.6 & 84.1 & \textbf{84.9} & 74.3 & 10.6 & 80.7 \\
XPaste & 86.2 & \underline{86.9} & \underline{$-0.7$} & 87.6 & 83.1 & 80.4 & \underline{2.8} & \underline{88.0} \\
\rowcolor{oursrow} \textbf{XPaste + AAS \emph{(ours)}} & \textbf{87.4} & \textbf{89.1} & \textbf{\boldmath $-1.8$} & \textbf{89.4} & 84.1 & \textbf{82.3} & \textbf{1.8} & \textbf{89.2} \\
SCP & 85.5 & 82.2 & 3.3 & 84.9 & 83.6 & 79.6 & 4.1 & 83.6 \\
\rowcolor{oursrow} SCP + AAS \emph{(ours)} & 86.9 & 84.9 & 2.0 & \underline{87.9} & \underline{84.4} & \underline{81.1} & 3.3 & 86.9 \\
\midrule
\multicolumn{9}{l}{\emph{Grounding DINO Swin-B (AP; paste rows: mini-COCO late-stage protocol)}} \\
Zero Shot & 94.1 & 90.5 & 3.6 & 93.0 & 90.3 & 88.5 & \underline{1.8} & 91.8 \\
XPaste & \underline{94.8} & \underline{92.6} & \underline{2.3} & 94.8 & \textbf{92.1} & 89.9 & 2.2 & \underline{93.8} \\
\rowcolor{oursrow} \textbf{XPaste + AAS \emph{(ours)}} & \textbf{95.3} & \textbf{93.8} & \textbf{1.5} & \textbf{95.8} & \underline{92.0} & \textbf{91.1} & \textbf{0.9} & \textbf{95.0} \\
SCP & 93.8 & 90.7 & 3.1 & 93.8 & 90.5 & 87.6 & 2.9 & 91.7 \\
\rowcolor{oursrow} SCP + AAS \emph{(ours)} & 94.6 & 92.2 & 2.4 & \underline{95.0} & \textbf{92.1} & \underline{90.3} & \underline{1.8} & 93.4 \\
\bottomrule
\end{tabular}

\end{table*}
